\documentclass[11pt]{article}

\usepackage[T1]{fontenc}
\usepackage[utf8]{inputenc}
\usepackage{amsmath,amssymb,amsthm}
\usepackage{mathtools}
\usepackage{mathrsfs}
\usepackage{booktabs}
\usepackage{geometry}
\usepackage[hidelinks]{hyperref}

\geometry{margin=1in}

\newtheorem{theorem}{Theorem}[section]
\newtheorem{lemma}[theorem]{Lemma}
\newtheorem{proposition}[theorem]{Proposition}
\newtheorem{corollary}[theorem]{Corollary}
\newtheorem{assumption}[theorem]{Assumption}
\theoremstyle{remark}
\newtheorem{remark}[theorem]{Remark}

\newcommand{\R}{\mathcal{R}}
\newcommand{\wideQ}{\widehat{Q}}
\newcommand{\sinc}{\operatorname{sinc}}

\title{A Scaling--Cancellation Principle for Chord-Based Geometric Approximations}
\author{C. Wayne Baker\\Independent Researcher\\Newfoundland and Labrador, Canada}
\date{Revised September 11, 2026}

\begin{document}

\maketitle

\begin{abstract}
This paper isolates a fixed-scale error-cancellation mechanism for geometric
approximations whose discretization error admits an even-power asymptotic
expansion.  If
\[
Q-Q_N
=
C_2N^{-2}+C_4N^{-4}+C_6N^{-6}+\cdots,
\]
then sampling the same approximation family at scaled resolutions gives
predictable factors $k^{-2j}$ on the successive error terms.  The two-scale
combination
\[
\wideQ_{N,k}
=
\frac{k^2Q_{kN}-Q_N}{k^2-1}
\]
therefore cancels the $N^{-2}$ term and satisfies
\[
Q-\wideQ_{N,k}
=
-\frac{C_4}{k^2}N^{-4}+O(N^{-6}).
\]
We then formulate the full multi-scale moment system, give its unique
Lagrange weights, and derive an explicit leading-error coefficient.  For
geometric scale nodes $N,bN,\ldots,b^sN$, cancellation of the first $s$
even-power terms yields
\[
Q-\wideQ_N^{(s)}
=
(-1)^s C_{2s+2}\,b^{-s(s+1)}N^{-2s-2}
+O(N^{-2s-4}).
\]
An equivalent recursive scale ladder gives a geometric analogue of the
Richardson--Romberg extrapolation table.  Applied to the half-perimeter of a
regular inscribed polygon,
$Q_N=N\sin(\pi/N)$, the tripling correction gives
\[
\pi-\frac{9Q_{3N}-Q_N}{8}
=
\frac{\pi^5}{1080}N^{-4}+O(N^{-6}).
\]
The framework also clarifies the distinction between fixed-resolution
scaling--cancellation and nonlinear residual-map refinement.  In particular,
the proportional-subtended angle-division residual
\[
\mathscr{R}_{a,b}(e)
=
-\frac{a^2-b^2}{24a^2}e^3+O(e^5)
\]
acts in the signed error variable and multiplies local error order under
iteration; it is not an instance of the fixed-scale extrapolation theorem.
No claim of new Richardson or Romberg theory is made.  The contribution is a
geometric formulation, an explicit fixed-scale hierarchy, and a careful
separation of the two convergence mechanisms.
\end{abstract}

\noindent\textbf{Keywords:}
Richardson extrapolation; geometric approximation; chord approximation;
polygonal approximation; error cancellation; N-series; asymptotic expansion.

\section{Introduction}

Many geometric quantities are approximated by replacing a smooth object with
a finite chordal, polygonal, or sampled object.  A resolution parameter $N$
then controls the approximation, and in favorable symmetric settings the
error has an expansion in even powers of $N^{-1}$:
\begin{equation}
\label{eq:even-expansion-intro}
Q-Q_N
=
C_2N^{-2}+C_4N^{-4}+C_6N^{-6}+\cdots.
\end{equation}
The essential observation is elementary: replacing $N$ by $kN$ multiplies
the term $N^{-2j}$ by $k^{-2j}$.  Once those scaling factors are known,
linear combinations of several resolutions can annihilate selected terms of
the expansion.

This is a geometric use of Richardson extrapolation, not a replacement for
it.  The point of the present paper is to state the mechanism in a form that
is directly reusable across a family of chord-based approximation problems
and to separate it cleanly from a second mechanism that occurs in
proportional-subtended angle refinement.

The paper makes four concrete steps.
\begin{enumerate}
\item It proves the two-scale cancellation formula and records the surviving
leading coefficient.
\item It gives the unique multi-scale weights through a Lagrange interpolation
formula in the scale variable $x=k^{-2}$.
\item It specializes the construction to geometric nodes
$N,bN,\ldots,b^sN$, obtaining the explicit factor $b^{-s(s+1)}$ in the
first surviving term and an equivalent recursive scale ladder.
\item It compares this fixed-resolution mechanism with a nonlinear cubic
angle-residual law, where the active variable is the signed error itself
rather than $N^{-2}$.
\end{enumerate}

The resulting framework is useful for polygonal approximations of $\pi$,
for chordal perimeter approximations when the required expansion has been
proved, and for the geometric N-series constructions developed in companion
work.  The same algebra also appears in classical extrapolation methods; the
present contribution is the geometric organization and the explicit
separation of convergence classes.

\section{Asymptotic hypothesis and scope}

The cancellation statements require an asymptotic expansion.  They do not
assert that every chord-based approximation possesses one.

\begin{assumption}[Even-power expansion through order $2r+2$]
\label{ass:even}
For some integer $r\ge 1$, the approximation family $Q_N$ satisfies
\begin{equation}
\label{eq:general-expansion}
Q-Q_N
=
\sum_{\ell=1}^{r+1} C_{2\ell}N^{-2\ell}
+O(N^{-2r-4})
\qquad (N\to\infty),
\end{equation}
where the coefficients $C_{2\ell}$ are independent of $N$.
\end{assumption}

In particular applications, an expansion of the form
\eqref{eq:general-expansion} may follow from a Taylor expansion, a periodic
quadrature argument, Euler--Maclaurin analysis, or a problem-specific exact
factorization.  Local symmetry often helps produce even powers, but symmetry
alone is not taken here as a universal theorem.

\begin{remark}[Admissible scale factors]
Because the index $N$ is typically integral, we use fixed integer scale
factors $k\ge2$ in the basic theorem.  More generally, the same algebra works
for any fixed scale factors for which all sampled resolutions are admissible.
\end{remark}

\begin{remark}[Claim boundary]
The results below are conditional extrapolation statements.  They do not
supply the expansion \eqref{eq:general-expansion} for an arbitrary curve,
mesh, parameterization, or chord rule.  Establishing that expansion remains
an application-specific analytic task.
\end{remark}

\section{The two-scale scaling--cancellation principle}

We begin with the elementary scaling law.

\begin{proposition}[Scaling of the asymptotic coefficients]
\label{prop:scaling}
Suppose
\[
Q-Q_N
=
C_2N^{-2}+C_4N^{-4}+C_6N^{-6}+O(N^{-8}).
\]
For a fixed integer $k\ge2$,
\[
Q-Q_{kN}
=
\frac{C_2}{k^2}N^{-2}
+
\frac{C_4}{k^4}N^{-4}
+
\frac{C_6}{k^6}N^{-6}
+O(N^{-8}).
\]
\end{proposition}

\begin{proof}
Substitute $kN$ for $N$ in the assumed expansion and use the fact that $k$
is fixed.
\end{proof}

\begin{theorem}[Two-scale scaling--cancellation]
\label{thm:two-scale}
Suppose
\begin{equation}
\label{eq:two-scale-assumption}
Q-Q_N
=
C_2N^{-2}+C_4N^{-4}+O(N^{-6}).
\end{equation}
For fixed $k\ge2$, define
\begin{equation}
\label{eq:two-scale-def}
\wideQ_{N,k}
=
Q_{kN}+\frac{Q_{kN}-Q_N}{k^2-1}
=
\frac{k^2Q_{kN}-Q_N}{k^2-1}.
\end{equation}
Then
\begin{equation}
\label{eq:two-scale-error}
Q-\wideQ_{N,k}
=
-\frac{C_4}{k^2}N^{-4}+O(N^{-6}).
\end{equation}
In particular, the order improves from $O(N^{-2})$ to $O(N^{-4})$.
\end{theorem}

\begin{proof}
From \eqref{eq:two-scale-assumption},
\[
Q_N
=
Q-C_2N^{-2}-C_4N^{-4}+O(N^{-6})
\]
and
\[
Q_{kN}
=
Q-\frac{C_2}{k^2}N^{-2}
-\frac{C_4}{k^4}N^{-4}+O(N^{-6}).
\]
Substituting these expressions into \eqref{eq:two-scale-def} gives
\[
\wideQ_{N,k}-Q
=
\frac{C_4}{k^2}N^{-4}+O(N^{-6}),
\]
which is equivalent to \eqref{eq:two-scale-error}.
\end{proof}

\begin{corollary}[Tripling]
\label{cor:tripling}
For $k=3$,
\begin{equation}
\label{eq:tripling}
\wideQ_{N,3}
=
\frac{9Q_{3N}-Q_N}{8},
\qquad
Q-\wideQ_{N,3}
=
-\frac{C_4}{9}N^{-4}+O(N^{-6}).
\end{equation}
\end{corollary}

The important point is that $C_2$ need not be known.  Its value disappears
because its scale factor $k^{-2}$ is known.

\section{Multi-scale cancellation as interpolation in the scale variable}

The two-scale formula extends naturally to several resolutions.  Let
\[
k_0=1,\qquad k_1,\ldots,k_s>0
\]
be distinct admissible scale factors and define
\[
x_j=k_j^{-2}.
\]
The asymptotic error at resolution $k_jN$ is a power series in $x_j$:
\[
Q-Q_{k_jN}
=
\sum_{\ell\ge1}C_{2\ell}x_j^{\ell}N^{-2\ell}.
\]
Thus cancellation of successive even powers is exactly a polynomial
interpolation condition at $x=0$.

\begin{theorem}[Unique $s$-term cancellation weights]
\label{thm:multiscale}
Assume
\begin{equation}
\label{eq:multi-expansion}
Q-Q_N
=
\sum_{\ell=1}^{s+1}C_{2\ell}N^{-2\ell}
+O(N^{-2s-4}).
\end{equation}
Let $k_0,\ldots,k_s$ be distinct positive admissible scale factors and set
$x_j=k_j^{-2}$.  Define
\begin{equation}
\label{eq:weights}
w_j
=
\prod_{\substack{m=0\\m\ne j}}^{s}
\frac{x_m}{x_m-x_j},
\qquad j=0,\ldots,s,
\end{equation}
and
\begin{equation}
\label{eq:multi-estimator}
\wideQ_N^{(s)}
=
\sum_{j=0}^{s}w_jQ_{k_jN}.
\end{equation}
Then
\begin{align}
\sum_{j=0}^{s}w_j&=1,\label{eq:mom0}\\
\sum_{j=0}^{s}w_jk_j^{-2\ell}&=0,
\qquad \ell=1,\ldots,s,\label{eq:moments}
\end{align}
and these conditions determine the weights uniquely.  Moreover,
\begin{equation}
\label{eq:leading-multiscale}
Q-\wideQ_N^{(s)}
=
(-1)^s C_{2s+2}
\left(\prod_{j=0}^{s}k_j^{-2}\right)
N^{-2s-2}
+O(N^{-2s-4}).
\end{equation}
\end{theorem}

\begin{proof}
The numbers $w_j$ in \eqref{eq:weights} are the Lagrange interpolation
weights for evaluating at $x=0$ from the distinct nodes
$x_0,\ldots,x_s$.  Therefore they reproduce constants and annihilate the
monomials $x,x^2,\ldots,x^s$, which proves
\eqref{eq:mom0}--\eqref{eq:moments}.  The associated Vandermonde system is
nonsingular, so the weights are unique.

Using \eqref{eq:multi-expansion} at each scale and
\eqref{eq:mom0}--\eqref{eq:moments}, all terms through $N^{-2s}$ cancel.
The first surviving coefficient is
\[
C_{2s+2}N^{-2s-2}
\sum_{j=0}^{s}w_jx_j^{s+1}.
\]
For $f(x)=x^{s+1}$, the interpolation error identity at $x=0$ gives
\[
0-\sum_{j=0}^{s}w_jx_j^{s+1}
=
\prod_{j=0}^{s}(0-x_j)
=
(-1)^{s+1}\prod_{j=0}^{s}x_j.
\]
Hence
\[
\sum_{j=0}^{s}w_jx_j^{s+1}
=
(-1)^s\prod_{j=0}^{s}x_j,
\]
which proves \eqref{eq:leading-multiscale}.
\end{proof}

\begin{remark}
The theorem is a Richardson extrapolation statement written in a geometric
scale variable.  The explicit product in \eqref{eq:leading-multiscale} is
useful because it makes the surviving coefficient transparent once the scale
nodes are fixed.
\end{remark}

\section{Geometric scale nodes and the N-series hierarchy}

A particularly natural choice in geometric refinement is a fixed base
$b\ge2$ with
\begin{equation}
\label{eq:geometric-nodes}
k_j=b^j,
\qquad
j=0,1,\ldots,s.
\end{equation}
These are the resolutions
\[
N,\ bN,\ b^2N,\ldots,b^sN.
\]
In the geometric N-series terminology used in companion work, these are the
successive scale samples.

\begin{corollary}[Geometric-node leading coefficient]
\label{cor:geometric-nodes}
Under the hypotheses of Theorem~\ref{thm:multiscale}, choose
$k_j=b^j$.  Then
\begin{equation}
\label{eq:geometric-leading}
Q-\wideQ_N^{(s)}
=
(-1)^s C_{2s+2}\,b^{-s(s+1)}N^{-2s-2}
+O(N^{-2s-4}).
\end{equation}
\end{corollary}

\begin{proof}
Since
\[
\prod_{j=0}^{s}k_j^{-2}
=
\prod_{j=0}^{s}b^{-2j}
=
b^{-2\sum_{j=0}^{s}j}
=
b^{-s(s+1)},
\]
the result follows from \eqref{eq:leading-multiscale}.
\end{proof}

The same hierarchy can be generated recursively, without solving the full
Vandermonde system at every order.

\begin{theorem}[Recursive fixed-base cancellation ladder]
\label{thm:recursive}
Let $Q_N^{[0]}=Q_N$.  For a fixed integer $b\ge2$, define recursively
\begin{equation}
\label{eq:recursive}
Q_N^{[r]}
=
\frac{b^{2r}Q_{bN}^{[r-1]}-Q_N^{[r-1]}}
{b^{2r}-1},
\qquad r\ge1.
\end{equation}
If
\[
Q-Q_N
=
\sum_{\ell=1}^{r+1}C_{2\ell}N^{-2\ell}
+O(N^{-2r-4}),
\]
then
\[
Q-Q_N^{[r]}
=
O(N^{-2r-2}).
\]
The value $Q_N^{[r]}$ is a linear combination of the $r+1$ geometric-scale
samples
\[
Q_N,Q_{bN},\ldots,Q_{b^rN}
\]
and agrees with the unique estimator of
Corollary~\ref{cor:geometric-nodes}.
\end{theorem}

\begin{proof}
Assume inductively that
\[
Q-Q_N^{[r-1]}
=
D_{2r}N^{-2r}+O(N^{-2r-2}).
\]
Then
\[
Q-Q_{bN}^{[r-1]}
=
D_{2r}b^{-2r}N^{-2r}+O(N^{-2r-2}).
\]
Substitution into \eqref{eq:recursive} cancels the $N^{-2r}$ term, leaving
$O(N^{-2r-2})$.  The recursion uses exactly the nodes
$N,bN,\ldots,b^rN$ and satisfies the same moment conditions as the unique
Lagrange estimator, so the two forms agree.
\end{proof}

\section{Regular polygon approximation of \texorpdfstring{$\pi$}{pi}}

For the unit circle, let
\begin{equation}
\label{eq:polygon-pi}
P_N
=
N\sin\!\left(\frac{\pi}{N}\right).
\end{equation}
This is half the perimeter of the regular inscribed $N$-gon, so
$P_N\to\pi$.  The sine expansion gives
\begin{equation}
\label{eq:pi-expansion}
\pi-P_N
=
\frac{\pi^3}{6}N^{-2}
-
\frac{\pi^5}{120}N^{-4}
+
\frac{\pi^7}{5040}N^{-6}
-
\frac{\pi^9}{362880}N^{-8}
+O(N^{-10}).
\end{equation}

Applying Corollary~\ref{cor:tripling} gives the one-step tripling estimator
\begin{equation}
\label{eq:pi-tripling}
\widehat P_N^{[1]}
=
\frac{9P_{3N}-P_N}{8}
\end{equation}
and the explicit error law
\begin{equation}
\label{eq:pi-tripling-error}
\pi-\widehat P_N^{[1]}
=
\frac{\pi^5}{1080}N^{-4}
+O(N^{-6}).
\end{equation}
Thus the raw $N^{-2}$ polygon error is converted into an $N^{-4}$ error
without using the value of the leading coefficient $\pi^3/6$ in the
extrapolation formula itself.

At the next level, the three tripling nodes $N,3N,9N$ give
\begin{equation}
\label{eq:pi-tripling-second}
\widehat P_N^{[2]}
=
\frac{P_N-90P_{3N}+729P_{9N}}{640}.
\end{equation}
By Corollary~\ref{cor:geometric-nodes},
\begin{equation}
\label{eq:pi-tripling-second-error}
\pi-\widehat P_N^{[2]}
=
\frac{\pi^7}{3674160}N^{-6}
+O(N^{-8}).
\end{equation}
Equations \eqref{eq:pi-tripling-error} and
\eqref{eq:pi-tripling-second-error} are examples of the general
$2,4,6,\ldots$ order ladder.

\begin{remark}
These formulas do not define a new value of $\pi$.  They accelerate a known
convergent polygonal approximation by cancelling known asymptotic modes.
\end{remark}

\section{Ellipse and smooth-curve applications}

The abstract theorem applies to any chordal perimeter family after an
expansion of the required form has been established for that family.  One
important companion example is the equal-parameter inscribed polygon of a
nondegenerate ellipse.  In that setting, the companion ellipse analysis gives
an exact factorization of the form
\begin{equation}
\label{eq:ellipse-factorization}
P_N
=
\sinc\!\left(\frac{\pi}{N}\right)M_N,
\qquad
\sinc x=\frac{\sin x}{x},
\end{equation}
where $M_N$ is a periodic midpoint approximation to the ellipse speed.  For a
fixed nondegenerate ellipse, the midpoint remainder is exponentially small,
so the algebraic part of the polygon error is governed by the sinc series.
Consequently,
\begin{equation}
\label{eq:ellipse-expansion}
\frac{P-P_N}{P}
=
\frac{\pi^2}{6N^2}
-
\frac{\pi^4}{120N^4}
+
O(N^{-6})+O(e^{-\rho N})
\end{equation}
for any admissible strip parameter $\rho>0$ below the analytic width.
The tripling correction therefore gives
\begin{equation}
\label{eq:ellipse-tripling}
P-\frac{9P_{3N}-P_N}{8}
=
\frac{P\pi^4}{1080}N^{-4}
+O(N^{-6})+O(e^{-\rho N}).
\end{equation}

For a more general smooth closed curve, one should not assume
\eqref{eq:ellipse-expansion}.  Instead, if a particular sampling and chord
rule is proved to satisfy
\[
P-P_N
=
C_2N^{-2}+C_4N^{-4}+O(N^{-6}),
\]
then Theorem~\ref{thm:two-scale} applies immediately.  This conditional
formulation keeps the extrapolation theorem separate from the analytic work
needed to justify the expansion in each geometry.

\section{A distinct mechanism: nonlinear cubic residual refinement}

The fixed-scale cancellation above acts in the resolution variable
$N^{-2}$.  A different error-suppression mechanism appears in the
proportional-subtended angle-division construction.

For positive integers $a$ and $b$, define the local transfer map
\[
\Phi_{a,b}(x)
=
2\arcsin\!\left(\frac{a}{b}\sin\frac{x}{2}\right).
\]
If $e$ is a signed angular error, the distinguished residual correction is
\begin{equation}
\label{eq:cubic-residual}
\mathscr{R}_{a,b}(e)
=
e-\Phi_{a,b}\!\left(\frac{b}{a}e\right).
\end{equation}
Its local expansion is
\begin{equation}
\label{eq:cubic-law}
\mathscr{R}_{a,b}(e)
=
-\frac{a^2-b^2}{24a^2}e^3+O(e^5),
\qquad a\ne b.
\end{equation}
For the motivating $4{:}3$ branch,
\begin{equation}
\label{eq:four-three}
\mathscr{R}_{4,3}(e)
=
-\frac{7}{384}e^3+O(e^5).
\end{equation}

The distinction from scaling--cancellation is structural.  If
$e=O(\varepsilon^p)$ and the cubic coefficient in
\eqref{eq:cubic-law} is nonzero, then
\[
\mathscr{R}_{a,b}(e)=O(\varepsilon^{3p}).
\]
Thus repeated residual correction multiplies the local error order by three:
\[
3\longrightarrow9\longrightarrow27\longrightarrow81\longrightarrow\cdots.
\]
By contrast, the fixed-scale hierarchy cancels successive asymptotic powers
of $N^{-2}$:
\[
2\longrightarrow4\longrightarrow6\longrightarrow8\longrightarrow\cdots.
\]

\begin{table}[ht]
\centering
\caption{Two distinct geometric error-suppression mechanisms.}
\label{tab:comparison}
\begin{tabular}{@{}lll@{}}
\toprule
Feature & Scaling--cancellation & Cubic residual refinement \\
\midrule
Active variable & $N^{-2}$ & signed error $e$ \\
Operation & linear combination of scales & nonlinear residual map \\
Local law & cancel selected powers & $e\mapsto ce^3+O(e^5)$ \\
Order behavior & $2,4,6,8,\ldots$ & $3,9,27,81,\ldots$ \\
Classical analogue & Richardson--Romberg & nonlinear local iteration \\
\bottomrule
\end{tabular}
\end{table}

The two mechanisms can coexist within a broader geometric approximation
program, but neither should be identified with the other.

\section{Classical extrapolation and claim boundary}

The algebra of Sections 3--5 belongs to the classical extrapolation family.
Richardson's work established the general principle of using asymptotic error
structure to improve approximations, and Romberg later organized repeated
extrapolation on geometrically refined quadrature meshes.  The recursion
\eqref{eq:recursive} is therefore not claimed as a new numerical-analysis
principle.

What is specific here is the geometric packaging:
\begin{itemize}
\item the resolution variable is attached directly to chordal and polygonal
geometric approximants;
\item the scale variable $x=k^{-2}$ makes the moment structure explicit;
\item geometric nodes yield the closed surviving factor
$b^{-s(s+1)}$;
\item the framework is placed beside, and explicitly separated from, the
nonlinear proportional-subtended cubic residual cascade.
\end{itemize}

Likewise, the phrase \emph{geometric N-series} is used here as terminology for
the author's scale-indexed geometric approximation framework.  It is not a
claim that the finite extrapolation formulas are outside classical
Richardson-type theory, nor is the fixed-scale mechanism identified with
Landen or modular transformations.

\section{Conclusion}

An even-power geometric error expansion provides a simple scale algebra.  If
\[
Q-Q_N
=
C_2N^{-2}+C_4N^{-4}+\cdots,
\]
then the samples $Q_{kN}$ carry the same coefficients multiplied by the known
factors $k^{-2j}$.  The two-scale estimator
\[
\wideQ_{N,k}
=
\frac{k^2Q_{kN}-Q_N}{k^2-1}
\]
cancels the leading term.  With $s+1$ distinct scales, Lagrange weights in
$x=k^{-2}$ cancel the first $s$ terms, and geometric scale nodes
$N,bN,\ldots,b^sN$ yield
\[
Q-\wideQ_N^{(s)}
=
(-1)^sC_{2s+2}b^{-s(s+1)}N^{-2s-2}
+O(N^{-2s-4}).
\]

For the regular polygon half-perimeter this gives a direct hierarchy of
accelerated approximations to $\pi$; for ellipse and other chordal perimeter
problems it applies once the needed expansion has been established.  The same
framework also provides a clean taxonomic boundary: fixed-resolution
Richardson-type cancellation raises the order along an even-power ladder,
whereas proportional-subtended residual refinement acts nonlinearly on the
signed error and multiplies its local order.

Future work may study stability of the multi-scale weights, optimal scale
selection under finite precision or geometric construction cost, and
extensions to curved-surface approximation families for which an even-power
expansion can be proved.

\begin{thebibliography}{9}

\bibitem{Richardson1911}
L.~F. Richardson,
``The approximate arithmetical solution by finite differences of physical
problems involving differential equations, with an application to the
stresses in a masonry dam,''
\emph{Philosophical Transactions of the Royal Society of London, Series A},
210 (1911), 307--357.

\bibitem{Romberg1955}
W.~Romberg,
``Vereinfachte numerische Integration,''
\emph{Norske Vid. Selsk. Forh., Trondheim},
28 (1955), 30--36.

\bibitem{BakerNSeries2026}
C.~W. Baker,
\emph{A Constructive N-Series Acceleration Law for Polygonal Approximation of Pi},
research manuscript and reproducibility archive, 2026.

\bibitem{BakerScaleOptimized2026}
C.~W. Baker,
\emph{Scale-Optimized Geometric Acceleration of Polygonal Approximations to Pi},
research manuscript and reproducibility archive, 2026.

\bibitem{BakerEllipseScaling2026}
C.~W. Baker,
\emph{A Universal Scaling Law for Ellipse Perimeters},
research manuscript and reproducibility archive, 2026.

\bibitem{BakerCubicRefinement2026}
C.~W. Baker,
\emph{A Local Cubic Refinement Law for Proportional-Subtended Angle Division},
research manuscript and reproducibility archive, 2026.

\end{thebibliography}

\end{document}
